\section{Datalog queries and rules}
\label{sec:datalog}

The definition of queries in ConceptBase is often complicated by the
limitations of the expressiveness of the query language or by the
limitations of the query optimizer to find the best solution.
The concept of datalog queries and rules was introduced to overcome
these limitations. Datalog queries and rules give the experienced user
the possibility to define the executable code of query (or rule) directly,
including the use of standard PROLOG predicates such as ground/1 to improve
the performance of a query or rule.

Although datalog queries and rules can be used as any other query (or rule),
they can cause an inconsistent database. This is due to the fact, that the
datalog queries and rules will usually not be evaluated while the semantic integrity
of the database is checked.

\subsection{Extended query model}

Datalog queries are defined in a similar way as standard query
classes. They must be declared as instance of the class {\ttfamily DatalogQueryClass}.

\begin{verbatim}
Class DatalogQueryClass isA GenericQueryClass with
attribute
   code : String
end
\end{verbatim}

The attribute {\ttfamily code} defines the executable code of the query as string.

Datalog rules have to be defined as an instance of {\ttfamily DatalogInRule} or {\ttfamily DatalogAttrRule},
depending on whether their conclusion should be an In-predicate or an A-predicate.

\begin{verbatim}
Class DatalogRule with
attribute
  concernedClass : Proposition;
  code : String
end

Class DatalogInRule isA DatalogRule
end

Class DatalogAttrRule isA DatalogRule
end
\end{verbatim}

The attribute {\ttfamily concernedClass} specifies the class for the In-predicate or the
attribute class for the A-predicate.

\subsection{Datalog code}

The datalog code is a list of predicates, separated by commas (,). As in Datalog or Prolog,
this will be interpreted as a conjunction of the predicates. To use disjunction, the {\ttfamily code}
attribute has to be specified multiple times.

All predicates that may be used in standard rules and queries may also be
used in datalog queries (see section \ref{sec:CBL} for a list). An argument of a predicate
may be one the following:

\begin{itemize}
\item an object name: If the object name starts with an upper case letter or includes special characters such as \verb+!+ or \verb+"+,
it must be written in single quotes (\verb+'+). This also holds for string object, e.g. \verb+"a string"+
must be written as \verb+'"a string"'+.
\item a predefined variable, defined by the context of the query or rule: If the query has a super class,
the variable {\ttfamily this} refers to the instances of this class. If the query has parameters
or computed attributes, variables with the names of the parameters or computed attributes will be predefined.
In a {\ttfamily DatalogInRule}, the variable {\ttfamily this} refers to the instances of the concerned class.
In a {\ttfamily DatalogAttrRule}, the variables {\ttfamily src} and {\ttfamily dst} refer to the source
and destination object of the attribute. Note, that {\bfseries all predefined variables have to be prefixed with
\verb+~+ and must be encoded in single quotes (\verb+'+)}, e.g. \verb+'~this'+, \verb+'~src'+, \verb+'~param'+.
\item existential variables: These variables must be declared in a special predicate {\ttfamily vars(list)},
which has to be the first predicate of the code. For example, the predicate \verb+vars([x,y])+ defines
the variables {\ttfamily x} and {\ttfamily y}.
\end{itemize}

A query expression of the form {\ttfamily query(q)} may be also used as predicate, or as second
argument of an In-predicate. {\ttfamily q} may be any valid query expression, e.g. just the
name of a query class, or a derive expression including the specification of parameters
(for example, \verb+find_instances[Class/class]+).

In addition, PROLOG predicates can be used as predicates. You can define your own PROLOG
predicates in a LPI-file (see section \ref{subsec:counter} for an example).

\subsection{Examples}

This section defines a few datalog queries and rules for the standard example model
of Employees, Departments and Managers (see \verb+$CB_HOME/examples/QUERIES+).

The first example defines a more efficient version of the recursive {\ttfamily MetabossQuery}.

\begin{verbatim}
DatalogQueryClass MetabossDatalogQuery isA Manager with
  parameter
     e : Employee
  code
     r1 : "In('~e','Employee'),A('~e',boss,'~this')";
     r2 : "vars([m]),
           In('~e','Employee'),
           In(m,query('MetabossDatalogQuery[~e/e]')),
           A(m,boss,'~this')"
end
\end{verbatim}

Note that the disjunction of the original query is represented by two code-attributes.
The example shows also the use of query expressions and existential variables.

The second example is the datalog version of the rule for the HighSalary class.
The infix-predicate \verb+>=+ is represented by the predicate {\ttfamily GE}.

\begin{verbatim}
Class HighSalary2 isA Integer
end

DatalogInRule HighSalaryRule2 with
concernedClass
   cc: HighSalary2
code
   c: "In('~this','Integer'),GE('~this',60000)"
end
\end{verbatim}


The last example shows the definition of a rule for an attribute. It also
shows, how the performance of a rule can be improved by specifying different
variants for different binding patterns. The example defines two rules,
depending on the binding of the variable {\ttfamily src}. The rule
defines the transitive closure of the {\ttfamily boss} attribute.
Rule {\ttfamily r1} is applied, if both arguments {\ttfamily src}
and {\ttfamily dst} are unbound. The second rule is used, if at least {\ttfamily src}
is bound, and the last rule will be applied, if we have a binding
for  {\ttfamily dst} but not for {\ttfamily src}.

\begin{verbatim}
DatalogAttrRule MetabossRule with
concernedClass
     cc : Employee!boss
code
    r1 : "vars([m]),var('~src'),var('~dst'),In(m,'Manager'),
          A(m,boss,'~dst'),A('~src',boss,m)";
    r2 : "vars([m]),ground('~src'),
          A('~src',boss,m),A(m,boss,'~dst')";
    r3 : "vars([m]),ground('~dst'),var('~src'),
          A(m,boss,'~dst'),A('~src',boss,m)"
end
\end{verbatim}

Note, that the predicates {\ttfamily var} and {\ttfamily ground}
are builtin predicates of PROLOG. Thus, this is also an example
for calling PROLOG predicates in a query or rule.
